% 原书第二篇第2.7节230--239；PDF页98--99，印刷页89--90。
% 本节凡碰撞频率z均指单分子每秒碰撞次数；
% 全系统碰撞事件数率记Z=(1/2)Nz。
\thermotransportproblem{230}{89}
理想气体压强在某过程中增大 $n$ 倍。
若过程分别为
(a) 等容；(b) 等温，
平均自由程 $\lambda$ 与每分子碰撞频率 $z$
各如何变化？
\exsolution
由硬球运动论
$\lambda\propto1/n_{\rm dens}$、
$z\propto n_{\rm dens}\sqrt T$。
(a) 等容时密度不变，
由 $p\propto T$ 得
$T_f/T_i=n$，
故
\[
\boxed{\lambda_f/\lambda_i=1,\qquad
z_f/z_i=\sqrt n}.
\]
(b) 等温时 $T$ 不变，
由 $p\propto n_{\rm dens}$ 得
$n_{{\rm dens},f}/n_{{\rm dens},i}=n$，
因此
\[
\boxed{\lambda_f/\lambda_i=1/n,\qquad
z_f/z_i=n}.
\]
\excheck 本题的 $n$ 是\emph{压强倍率}，
而 $n_{\rm dens}$ 表示分子数密度，
两者不能混为同一物理量。

\thermotransportproblem{231}{89}
由刚性双原子分子组成的理想气体
经历可逆绝热过程。
求平均自由程 $\lambda$
和每分子碰撞频率 $z$ 分别依赖于：
(a) 容积 $V$；(b) 压强 $p$；
(c) 绝对温度 $T$ 的幂律。
\exsolution
刚性双原子气体
$\gamma=7/5$。
分子数密度 $n_{\rm dens}\propto1/V$，
故 $\lambda\propto V$；
而
$z\propto n_{\rm dens}\sqrt T
\propto V^{-1}T^{1/2}$。
绝热关系
$TV^{\gamma-1}=\text{常数}$，
故 $T\propto V^{-2/5}$。
所以
\[
\boxed{\lambda\propto V,\qquad
z\propto V^{-6/5}}.
\]
又
$pV^\gamma=\text{常数}$，
$V\propto p^{-5/7}$，
得到
\[
\boxed{\lambda\propto p^{-5/7},\qquad
z\propto p^{6/7}}.
\]
最后
$V\propto T^{-5/2}$，
因此
\[
\boxed{\lambda\propto T^{-5/2},\qquad
z\propto T^{3}}.
\]
\excheck 最后一项须由
$V^{-6/5}\propto(T^{-5/2})^{-6/5}=T^3$
计算；若简答文本抽取显示 $T^2$，
与同题前两种表示不能同时相容，应依据连续性核验。

\thermotransportproblem{232}{89}
理想气体经历多方过程
$pV^n=\text{常数}$。
求平均自由程 $\lambda$
和每分子碰撞频率 $z$
分别随
(a) $V$；(b) $p$；(c) $T$
的幂律。
\exsolution
仍有
$\lambda\propto V$、
$z\propto V^{-1}\sqrt T$。
多方状态关系给
$T=pV/(\nu R)\propto V^{1-n}$，
所以
\[
\boxed{\lambda\propto V,\qquad
z\propto V^{-(n+1)/2}}.
\]
若 $n\ne0$，
$p\propto V^{-n}$，
$V\propto p^{-1/n}$，
故
\[
\boxed{\lambda\propto p^{-1/n},\qquad
z\propto p^{(n+1)/(2n)}}.
\]
若 $n\ne1$，
$V\propto T^{1/(1-n)}$，
得到
\[
\boxed{\lambda\propto T^{1/(1-n)},\qquad
z\propto T^{(n+1)/[2(n-1)]}}.
\]
\excheck 特例 $n=0$ 为等压过程，
不能以 $p$ 为随过程变化的自变量；
$n=1$ 是等温过程，
不能用 $T$ 的变化幂律描述。

\thermotransportproblem{233}{89}
刚性双原子理想气体经历多方过程。
若分子间碰撞事件的次数率保持恒定，
求过程摩尔热容：
(a) 每单位体积的碰撞事件率不变；
(b) 整团气体的碰撞事件总率不变。
\exsolution
记分子数密度为 $n_d=N/V$，
单分子碰撞频率
$z\propto n_d\sqrt T$。
\emph{每单位体积的碰撞事件率}
$Z_v=(1/2)n_dz
\propto n_d^2\sqrt T
\propto V^{-2}T^{1/2}$。
(a) 若 $Z_v$ 不变，
则 $T\propto V^4$。
由 $T\propto V^{1-n}$，
多方指数为 $n=-3$。
双原子气体 $C_V=5R/2$，
\[
\boxed{C_a=C_V+\frac R{1-n}
=\frac{5R}{2}+\frac R4
=\frac{11R}{4}\simeq23\,\mathrm{J/(mol\,K)}}.
\]
(b) 若总碰撞率
$Z_{\rm tot}=VZ_v\propto V^{-1}\sqrt T$
不变，则
$T\propto V^2$、$n=-1$，
\[
\boxed{C_b=C_V+\frac R{1-(-1)}
=3R\simeq25\,\mathrm{J/(mol\,K)}}.
\]
\excheck \textbf{与书末简答存在定义冲突：}
本版简答(b) 印 $C=7R/2\simeq29\,\mathrm{J/(mol\,K)}$，
但若“碰撞次数”严格指\emph{单位时间内的碰撞事件率}，
则 $Z_{\rm tot}\propto V^{-1}\sqrt T$ 恒定导出 $3R$。
已明确碰撞计数口径；
若底本采用其他时间尺度上的“碰撞次数”，
需进一步查证，暂不把两答案强行等同。

\thermotransportproblem{234}{89}
摩尔质量为 $M$ 的理想气体装在
体积 $V$ 的薄壁容器中，
容器壁维持温度 $T$。
在 $t=0$ 时打开面积很小的孔口 $S$，
气体开始向真空逸出。
若初始数密度为 $n_0$，
求数密度 $n(t)$。
\exsolution
容器壁保持温度，
可视为逸出气体内部分布始终处于
温度 $T$ 的准平衡麦克斯韦分布。
单位面积每秒外逸粒子通量
$\Phi=n\bar v/4$，
其中
$\bar v=\sqrt{8RT/(\pi M)}$。
容器粒子数 $N=nV$，故
\[
\frac{\mathrm dN}{\mathrm dt}
=-S\frac{n\bar v}{4}
=-S\frac{N\bar v}{4V}.
\]
因此
$\mathrm dN/\mathrm dt=-(1/\tau)N$，
时间常数
$\tau=4V/(S\bar v)$。
解得
\[
\boxed{n(t)=n_0e^{-t/\tau}},\qquad
\boxed{\tau=\frac{4V}{S\bar v}}.
\]
\excheck 需是孔径小于分子平均自由程尺度的
\emph{分子逸出}，不能直接套连续流中的喷嘴方程。

\thermotransportproblem{235}{89}
一容器被隔热薄板分成相等两腔。
隔板有两个孔：一个很小、
一个尺寸远大于分子平均自由程。
第二腔始终维持在第一腔温度的
$\eta$ 倍。
当关闭大孔、只让小孔连通后，
第二腔气体数密度变化几倍？
\exsolution
大孔打开时，
两腔可达常规流体力学的共同压强：
$p_1=p_2$。
由 $p_i=n_i k_BT_i$ 和
$T_2=\eta T_1$，
初态有
$n_{2,0}/n_{1,0}=1/\eta$。
关闭大孔而保留\emph{分子流}小孔时，
分子逸出平衡需要两方向通量相等：
\[
n_{1,f}\bar v_1=n_{2,f}\bar v_2.
\]
因为 $\bar v\propto\sqrt T$，
得到
$n_{2,f}/n_{1,f}=1/\sqrt\eta$。
两腔体积相等且总分子数不变，
$n_1+n_2$ 在前后相等，
于是
\[
n_{2,0}=\frac{n_1+n_2}{1+\eta},
\qquad
n_{2,f}=\frac{n_1+n_2}{1+\sqrt\eta}.
\]
故
\[
\boxed{\frac{n_{2,f}}{n_{2,0}}
=\frac{1+\eta}{1+\sqrt\eta}}.
\]
\excheck $\eta>1$ 时该比值大于1，
即高温第二腔分子浓度增加，
这是一种热渗透/分子逸出平衡效应。

\thermotransportproblem{236}{89}
理想气体经过某过程后，
动力黏度系数增大
$\alpha=2.0$ 倍，
扩散系数增大
$\beta=4.0$ 倍。
求压强变化的倍率。
\exsolution
稀薄气体运动论给
\[
\eta_{\rm visc}\propto\sqrt T,\qquad
D\propto\frac{\sqrt T}{n_d},
\]
其中 $n_d$ 是分子数密度。
黏度增至 $\alpha$ 倍意味着
$T_f/T_i=\alpha^2$。
扩散系数增至 $\beta$ 倍：
\[
\beta=\frac{D_f}{D_i}
=\alpha\frac{n_{d,i}}{n_{d,f}},
\]
故
$n_{d,f}/n_{d,i}=\alpha/\beta$。
由 $p=n_dk_BT$，
\[
\boxed{\frac{p_f}{p_i}
=\frac{\alpha}{\beta}\alpha^2
=\frac{\alpha^3}{\beta}
=\frac{2^3}{4}=2}.
\]
\excheck 原书答案行纯文本看似
“$\alpha^2/\beta$”，但从
$T\sim\eta_{\rm visc}^2$ 与 $D\sim\eta/n_d$
的两个关系检查，压力倍率应为
$\alpha^3/\beta$；给定数值恰为2，
与书末数值答案一致。

\thermotransportproblem{237}{89}
如果理想气体体积增大 $n$ 倍，
求扩散系数 $D$ 与动力黏度 $\eta$
的变化倍数：
(a) 等温膨胀；
(b) 等压膨胀。
\exsolution
动理论关系
$D\sim\bar v\lambda$、
$\eta\sim\rho\bar v\lambda$。
对同一种气体
$\bar v\propto T^{1/2}$，
$\lambda\propto1/n_d\propto V$，
而 $\rho\propto1/V$，因此
$\eta\propto T^{1/2}$、
$D\propto V T^{1/2}$。
(a) $T$ 不变时：
\[
\boxed{\frac{D_f}{D_i}=n,\qquad
\frac{\eta_f}{\eta_i}=1}.
\]
(b) $p$ 不变时理想气体
$T\propto V$，所以：
\[
\boxed{\frac{D_f}{D_i}=n^{3/2},\qquad
\frac{\eta_f}{\eta_i}=n^{1/2}}.
\]
\excheck 黏度不是一般地正比于气体密度；
在稀薄气体动理论中，密度和平均自由程
的变化彼此抵消，留下 $\sqrt T$ 因子。

\thermotransportproblem{238}{89}
刚性双原子理想气体可逆绝热压缩，
体积变成原来的 $1/n$，
其中 $n=10$。
求扩散系数 $D$ 与黏度 $\eta$
各增减多少倍。
\exsolution
刚性双原子气体
$\gamma=7/5$，
绝热压缩有
$T_f/T_i=(V_i/V_f)^{\gamma-1}
=n^{2/5}$。
平均速率变化
$\bar v_f/\bar v_i=n^{1/5}$。
平均自由程与体积成正比，
$\lambda_f/\lambda_i=1/n$，
因此
\[
\frac{D_f}{D_i}
=\frac{\bar v_f}{\bar v_i}
\frac{\lambda_f}{\lambda_i}
=n^{1/5-1}=n^{-4/5},
\]
即
\[
\boxed{D\text{ 减小 }n^{4/5}=10^{0.8}
\simeq6.3\text{ 倍}}.
\]
黏度只随 $\sqrt T$ 变化，
\[
\boxed{\eta\text{ 增大 }n^{1/5}
=10^{0.2}\simeq1.6\text{ 倍}}.
\]
\excheck 不可把扩散系数和黏度
都视为同一密度幂律；二者对体积的依赖不同。

\thermotransportproblem{239}{90}
理想气体沿多方过程
$pV^n=\text{常数}$。
如果在该过程中下列系数保持不变，
分别求多方指数 $n$：
(a) 扩散系数 $D$；
(b) 动力黏度 $\eta$；
(c) 热导率 $\kappa$。
\exsolution
多方关系给
$T\propto V^{1-n}$。
稀薄气体中
$D\propto V\sqrt T
\propto V^{(3-n)/2}$。
(a) 要使指数为零：
\[
\boxed{n_D=3}.
\]
动力黏度
$\eta\propto\sqrt T
\propto V^{(1-n)/2}$，
故
\[
\boxed{n_\eta=1}.
\]
对同一种气体、定容质量比热近似常数，
热导率
$\kappa\propto\rho\bar v\lambda c_V
\propto\sqrt T$，
所以
\[
\boxed{n_\kappa=1}.
\]
\excheck $n=1$ 是等温过程，
而 $n=3$ 是非等温的特殊多方过程；
气体的三种输运系数无需同时不变。
